\documentclass[12pt,a4paper]{article}
\usepackage[utf8]{inputenc}
\usepackage[T3,T1]{fontenc}
\usepackage[portuges]{babel}
\usepackage[noenc]{tipa}
\usepackage{tipx}
\usepackage[geometry,weather,misc,clock]{ifsym}
\usepackage{pifont}
\usepackage{eurosym}
\usepackage{amsmath}
\usepackage{wasysym}
\usepackage{amssymb,amsfonts,textcomp}
\usepackage{color}
\usepackage{array}
\usepackage{hhline}
\usepackage{booktabs}
\usepackage{hyperref}
\hypersetup{pdftex, colorlinks=true, linkcolor=blue, citecolor=blue,
            filecolor=blue, urlcolor=blue, pdftitle=}
\usepackage[pdftex]{graphicx}

% Page layout (geometry)
\setlength\voffset{-1in}
\setlength\hoffset{-1in}
\setlength\topmargin{0.7874in}
\setlength\oddsidemargin{0.9846in}
\setlength\textheight{10.118099in}
\setlength\textwidth{6.2988997in}
\setlength\footskip{0.0cm}
\setlength\headheight{0cm}
\setlength\headsep{0cm}

% Pages styles
\makeatletter
\newcommand\ps@Standard{
  \renewcommand\@oddhead{}
  \renewcommand\@evenhead{}
  \renewcommand\@oddfoot{}
  \renewcommand\@evenfoot{}
  \renewcommand\thepage{\arabic{page}}
}
\makeatother
\pagestyle{Standard}
\setlength\tabcolsep{1mm}
\renewcommand\arraystretch{1.3}


\begin{document}

\begin{tabular}{lcr}
    \includegraphics{img/logo-on.jpg}
    & \hspace{5cm} &
    \includegraphics{img/logo-mcti.png}
\end{tabular}

\vspace{0.5cm}

\begin{center}
    \textbf{\LARGE RELATÓRIO 2015}

    \vspace{1cm}

    \textbf{MODELAGEM E INVERSÃO DE CAMPOS GRAVITACIONAIS EM COORDENADAS
            ESFÉRICAS}
\end{center}

\vspace{1.5cm}

\begin{flushleft}

{\bfseries
NOME DO ALUNO:
}
Leonardo Uieda

{\bfseries
NIVEL:
}
Doutorado

{\bfseries
ENDEREÇO RESIDENCIAL:
}
Rua Aguiar, 49 ap 104. 20261-120 Tijuca, Rio de Janeiro - RJ

{\bfseries
TELEFONES FIXO e CELULAR:
}
(21) 983636761

{\bfseries
EMAIL NÃO INSTITUCIONAL:
}
leouieda@gmail.com

{\bfseries
DATA DE INÍCIO NA PÓS-GRADUAÇÃO:
}
11/2011

{\bfseries
DATA DE INÍCIO DO PROJETO DE PESQUISA:
}
11/2011

{\bfseries
NOME DO ORIENTADOR:
}
Valéria C. F. Barbosa

{\bfseries
PERÍODO PREVISTO DE BOLSA DE ESTUDOS:
}
11/2011 - 10/2015

{\bfseries
PERÍODO A QUE SE REFERE O RELATÓRIO:
}
03/2014 - 02/2015

\vfill

{\small Este relatório foi entregue ao orientador no dia}
\rule{1.5cm}{0.4pt}
\hfill
\rule{4cm}{0.4pt}

\vspace{-0.2cm}

\hfill {\tiny ASSINATURA DO ORIENTADOR}

\vspace{0.5cm}

{\small O orientador emitiu parecer no dia}
\rule{1.5cm}{0.4pt}
\hfill
\rule{4cm}{0.4pt}

\vspace{-0.2cm}
\hfill {\tiny ASSINATURA DO ORIENTADOR}

\vspace{0.5cm}

{\small Este relatório e o parecer do orientador foi recebido na
SECRETARIA DA PÓS-GRADUAÇÃO no dia:}

\vspace{0.5cm}

\rule{3cm}{0.4pt}
\hfill
\rule{10cm}{0.4pt}

\vspace{-0.2cm}

\hfill {\tiny NOME E ASSINATURA DO FUNCIONÁRIO DA SECRETARIA}

\end{flushleft}

\newpage

%-----------------------------------------------------------------------------
% PAGINA DE INFORMACOES CURRICULARES
%-----------------------------------------------------------------------------
\begin{center}
\textbf{\large INFORMAÇÕES CURRICULARES, CONTENDO OS SEGUINTES ITENS:}
\end{center}

\vspace{1cm}

\begin{flushleft}

\noindent (1) TOTAL DE CRÉDITOS CURSADOS EM DISCIPLINAS: \rule{2cm}{0.4pt}

\bigskip

\noindent (2) LISTA DE TODAS AS DISCIPLINAS CURSADAS ATÉ O MOMENTO PELO ALUNO
E OS CONCEITOS OBTIDOS

\bigskip

Física da Terra - A\\
Introdução à geofísica de exploração - A\\
Instrumentação geofísica - A\\
Problemas inversos em geofísica - A\\
Métodos Matemáticos - B\\
Métodos computacionais aplicados à geofísica - A\\
Problemas não lineares - Métodos por gradiente - A\\
Problemas não lineares - Métodos por heurísticos - A\\
Estado térmico da Terra - A\\
Tópicos de interpretação de dados gravimétricos e magnéticos - A

\bigskip

\noindent (3) SITUAÇÃO DO ALUNO QUANTO AOS CRÉDITOS OBTIDOS NOS SEMINÁRIOS
ANUAIS

\bigskip

Aprovado nos anos de 2012 e 2014.

\bigskip

\noindent (4) SITUAÇÃO DO ALUNO QUANTO AO EXAME DE PROFICIÊNCIA

\bigskip

Aprovado.

\bigskip

\noindent (5) SITUAÇÃO DO ALUNO DE DOUTORADO QUANTO AO EXAME DE QUALIFICAÇÃO

\bigskip

Aprovado no ano 2013.

\bigskip

\noindent (6) LISTA DAS REUNIÕES CIENTÍFICAS EM QUE PARTICIPOU NO PERÍODO A QUE
SE REFERE O RELATÓRIO, COM O TÍTULO E AUTORES DOS TRABALHOS APRESENTADOS

\bigskip

Evento: Python in Science Conference 2014\\
Autores: Uieda, L., V. C. Oliveira Jr and V. C. F. Barbosa\\
Título: Using Fatiando a Terra to solve inverse problems in geophysics\\
Material disponível em https://github.com/leouieda/scipy2014

\bigskip

Evento: EGU General Assembly 2014\\
Autores: Uieda, L. and V. C. F. Barbosa\\
Título: Gravity inversion in spherical coordinates using tesseroids\\
Material disponível em https://github.com/leouieda/egu2014

\bigskip

\noindent (7) LISTA DOS ARTIGOS PUBLICADOS, ACEITOS OU SUBMETIDOS

\bigskip

(Submetido) Carlos, D. U., Uieda, L., and V. C. F. Barbosa (2015),
How two gravity-gradient inversion methods can be used to reveal different
geologic features of ore deposit - a case study from the Quadrilátero
Ferrífero (Brazil),
Journal of Applied Geophysics.

\bigskip

(Em revisão) Oliveira Jr, V. C., D. P. Sales, V. C. F. Barbosa, and L. Uieda (2014),
Estimation of the total magnetization direction of approximately spherical
bodies, Nonlinear Processes in Geophysics.

\bigskip

(Publicado) Uieda, L., V. C. Oliveira Jr, and V. C. F. Barbosa (2014), Geophysical
tutorial: Euler deconvolution of potential-field data, The Leading Edge, 33(4),
448-450, doi:10.1190/tle33040448.1.

\bigskip

(Publicado) Carlos, D. U., L. Uieda, and V. C. F. Barbosa (2014), Imaging iron ore from the
Quadrilátero Ferrífero (Brazil) using geophysical inversion and drill hole
data, Ore Geology Reviews, 61, 268-285, doi:10.1016/j.oregeorev.2014.02.011.

\bigskip

\noindent (8) OUTRAS ATIVIDADES RELEVANTES NO PERÍODO (PARTICIPAÇÃO EM
TRABALHOS DE CAMPO, ESCOLAS ETC)

\bigskip

Ministrei o curso "Tópicos de inversão em geofísica" com duração de dois dias
na III Semana de Geofísica da UNB.
O material didático e informações sobre o curso estão disponíveis em
https://github.com/pinga-lab/inversao-unb-2014


\bigskip

\end{flushleft}


\newpage

%-----------------------------------------------------------------------------
% PROJETO
%-----------------------------------------------------------------------------
\begin{center}\textbf{\large PROJETO ORIGINAL}\end{center}

\vspace{1cm}

\textit{
Nesta parte do relatório o estudante deve incluir o projeto de pesquisa tal
como apresentado à Comissão de Pós-Graduação em Geofísica do Observatório
Nacional na época da inscrição.
}

\newpage

%-----------------------------------------------------------------------------
% DESENVOLVIMENTO DO PROJETO DE PESQUISA NO ÚLTIMO RELATÓRIO
%-----------------------------------------------------------------------------
\begin{center}
\textbf{\large RELATO DO DESENVOLVIMENTO DO PROJETO DE PESQUISA NO ÚLTIMO
RELATÓRIO}
\end{center}

\vspace{1cm}

No período entre Março de 2013 e Fevereiro de 2014,
eu desenvolvi a infraestrutura computacional que seria necessária para a
implementação da inversão de dados gravimétricos em coordenadas esférias.

A primeira etapa foi o desenvolvimento de um conjunto
de classes e funções em linguagem Python para
automatizar a criação de métodos de inversão.
Esse código foi incluído na biblioteca Fatiando a Terra (http://fatiando.org/),
onde é utilizado para implementar os diversos métodos de inversão já incluídos
na biblioteca.
Esse novo módulo de inversão automatiza as tarefas de
(1) construção da função objetivo,
(2) implementação das diversas formas de regularização,
(3) minimização da função objetivo por diferentes métodos,
(4) cálculo do vetor gradiente e matriz Hessiana da função objetivo
e (5) determinação do parâmetro de regularização através da análise da curva L.

A segunda etapa foi a implementação da modelagem direta em coordenadas
esféricas utilizando tesseroides (prismas esféricos).
Fiz uma implementação da modelagem em linguagem Python para comparar com a
implementação anterior em linguagem C do programa Tesseroids
(http://www.leouieda.com/tesseroids).
Obtive resultados preliminares da análise do erro cometido na integração
numérica do campo gravitacional de um tesseroide.
Existem poucos modelos em coordenadas esféricas com soluções analíticas que
possam ser utilizados na quantificação do erro.
Utilizei o modelo de meia casca esférica, para a qual existe uma solução
analítica para o ponto de observação localizado no polo.
Outro modelo utilizado foi o de um prisma retangular reto.
Este não é um modelo ideal pois a aproximação de um tesseroide por um prisma
retangular piora quanto maior o tesseroide ou quanto mais longe da região
equatorial.



\newpage

%-----------------------------------------------------------------------------
% TRABALHO DE PESQUISA DESENVOLVIDO NO PERÍODO
%-----------------------------------------------------------------------------
\begin{center}
\textbf{\large DESCRIÇÃO DETALHADA DO TRABALHO DE PESQUISA DESENVOLVIDO NO
PERÍODO DO RELATÓRIO}
\end{center}

\vspace{1cm}

\textit{
Esta parte do relatório deve conter \textbf{obrigatoriamente}
os seguintes itens:
}

\vspace{1cm}

{\centering\bfseries\itshape
Metodologia aplicada ou desenvolvida
\par}

\bigskip

{\centering\bfseries\itshape
Resultados parciais já obtidos
\par}

\bigskip

{\centering\bfseries\itshape
Dificuldades encontradas e como elas estão sendo superadas
\par}

\bigskip

{\centering\bfseries\itshape
Bibliografia utilizada no contexto do trabalho.
\par}

\bigskip

\bigskip

{\itshape
Recomenda-se que o estudante escreva, nesta parte do relatório, textos que
poderão ser aproveitados em suas teses ou dissertações.}

\bigskip

{\itshape
Por exemplo, aqueles alunos que realizaram trabalhos de campo ou de laboratório
no último ano e estes trabalhos fazem parte de seu projeto de pesquisa; então,
estes alunos podem aproveitar para escrever sobre estes trabalhos executados.}

\bigskip

{\itshape
Outro exemplo, aqueles alunos que estão desenvolvendo metodologias ou
aplicações de metodologias existente; então, estes alunos podem aproveitar
para escrever sobre estas metodologias.}

\bigskip

{\itshape
Outro exemplo, aqueles alunos mais adiantados que já escreveram nos relatórios
anteriores os itens acima; então, estes alunos podem aproveitar para escrever
sobre seus resultados apresentando figuras que serão aproveitadas na Tese ou
Dissertação.}

\bigskip

{\itshape
Alguns alunos também podem escrever o texto da introdução do seu trabalho. Em
geral, uma introdução deve: 1) apresentar a natureza e o alcance do problema;
2) revisar a literatura;, 3) apresentar os objetivos do trabalho; 4) descrever
o método de investigação; e 5) descrever os principais resultados da
investigação.}



\newpage

%-----------------------------------------------------------------------------
% PRÓXIMAS ETAPAS DO TRABALHO DE PESQUISA
%-----------------------------------------------------------------------------
\begin{center}
\textbf{\large PRÓXIMAS ETAPAS DO TRABALHO DE PESQUISA}
\end{center}

\vspace{0.5cm}

\begin{flushleft}

\textbf{Atividades de pesquisa previstas para o próximo período:}

\begin{enumerate}
    \item Finalização da escrita do artigo sobre a modelagem direta utilizando
        tesseroides. Submissão para a revista \textit{Geophysics}.
    \item Aplicação a dados sintéticos e dados reais da inversão para
        determinação do relevo de interfaces em coordenadas esféricas.
    \item Escrita e submissão do artigo sobre a inversão de interfaces.
    \item Elaboração da tese baseada nos dois artigos.
    \item Defesa.
\end{enumerate}

\textbf{Atividades acadêmicas previstas para o próximo período:}

\bigskip

Nenhuma.

\bigskip

\textbf{Cronograma detalhado das atividades:}

\begin{center}
\begin{tabular}{ll}
    \toprule
    \textbf{Atividade} & \textbf{Período de execução} \\
    \midrule
    Escrita do artigo de modelagem direta & Fev-Março\\
    Aplicação da inversão de interfaces & Março-Maio\\
    Escrita do artigo de inversão de interfaces & Abril-Agosto\\
    Finalização da tese & Agosto-Setembro\\
    Defesa & Outubro\\
    \bottomrule
\end{tabular}
\end{center}


\bigskip

\textbf{Data prevista de conclusão do mestrado ou doutorado:} Outubro/2015

\end{flushleft}


\end{document}
